\subsubsection{Invisible decays of Higgs and other mesons in models with singlet neutrinos in large extra dimensions --- arXiv:hep-ph/0303160}

\textbf{Authors:} N. G. Deshpande, Dilip Kumar Ghosh

\paragraph{主な主張}
Bulk singlet neutrino の多数の KK mode によって Higgs invisible decay は増大し得るが、ニュートリノ質量と高次元 Yukawa coupling の摂動性を両立させると、$H\to b\bar b$ に匹敵する幅が得られる領域は非常に限られることを示した \cite{Deshpande:2003Invisible}。

\paragraph{新規性}
KK mode の多重度による Higgs invisible width の増強を、当時の大気ニュートリノ振動データと Yukawa coupling の大きさに照らして定量評価し、右巻きニュートリノが重力より少ない次元を伝播する場合や非対称な余剰次元も検討した。

\paragraph{位置付け}
Bulk-neutrino 模型での Higgs invisible decay の初期の現象論研究であり、現在の $T^2$ 模型における KK 状態密度による崩壊幅の増強を比較する先行研究となるが、一般のトーラス複素構造への依存性は扱っていない。

\paragraph{コメント：Yukawa coupling の摂動性}
本文では高次元で定義された無次元結合 $\lambda$ が $\mathcal{O}(10\text{--}20)$ を超えると非摂動的になると見積もり、数値評価では保守的に $\lambda\leq 10$ を採用する。これは論文中でループ計算から厳密に導いた上限ではなく、素朴な4次元のループ展開パラメータ $\lambda^2/(16\pi^2)$ が小さくなくなることを踏まえた目安である。個々の KK mode への4次元有効結合は $Y_\nu\sim\lambda M_*/M_{\mathrm{Pl}}$ と小さいが、多数の KK mode の和が補正を増強し得るため、高次元理論全体の摂動性と区別する必要がある。論文では $\kappa\equiv\Gamma(H\to\nu\bar\nu)/\Gamma(H\to b\bar b)\sim 1$ の領域で $\lambda\sim80$ が必要となり、$\lambda\leq10$ を課すと $M_*$ の下限が高くなって $\kappa$ が抑制される；この数値基準を6次元 $T^2$ 模型へそのまま移すには、brane Yukawa coupling の正規化と KK 和を改めて確認する必要がある。
